% Incorporación conservativa: selección aditiva e instrumento térmico.
\section{Sélection additive et instrument de la lecture thermique}
\label{sec:ii-aditiva-instrumento}

La référence spectrale et l’action construites dans les sections précédentes
permettent de préciser ce qui est compté lorsqu’on mesure l’information et l’énergie. Le coefficient
thermique apparaît comme une intensité relative entre deux événements du même
registre. Sa fonction informationnelle est caractérisée par l’additivité, et
les deux événements admettent un appareil explicite qui conserve l’état complet.
On développe ainsi la règle marquée de
\ref{thm:ii-ter27-transducer}, avec sa réalisation, sa sélection au sein d’une
classe de lecteurs et ses conditions de transport.

La composition maintient l’antécédent APP--TRIT--TPK : feuillets additif et
multiplicatif, régime, orientation, résidu, quotient, frontière, chemin et
mémoire. Le prolongement
$w_6\to w_{12}\to w_{18}\to w_{24}\to w_{30}\to R_{36}\to G_9$
et l’ordre $U_t\to\Gamma_9\to K_{\rm ph}$ précèdent cette lecture
spectrale. Le retour de phase prolonge le registre ; les marques observées
ne remplacent ni ce registre ni les deux projections de la structure discrète
conjointe du continu. On reçoit la raison $r=1/10$, les rangs
$(1,380,649)$ et leur assignation aux degrés $23+3n$ avec les conditions de
\ref{subsec:ii-ter27-reference}. Cette assignation est conservée comme partie
explicite de la réalisation, avec le hamiltonien et son origine énergétique.

% BEGIN THERMAL_R03_01
\subsection{Événements spectraux et mesure relative}
\label{subsec:ii-tins-eventos}

Soit $\mathcal K=\ell^2(\mathbb N_0)\otimes\mathbb C^M$, avec $M=649$,
et considérons la densité normalisée
\begin{equation}
 \Sigma=\rho_r\otimes I_M/M,\qquad
 \rho_r=(1-r)\sum_{j\ge0}r^j|j\rangle\langle j|.
 \label{eq:ii-tins-sigma}
\end{equation}
Pour $n=0,1,2$, soient $j_n=23+3n$ et $Q_n$ des projecteurs de mémoire de rang
$d_n$, où $(d_0,d_1,d_2)=(1,380,649)$. Fixons un vecteur unitaire $\chi_B$.
Les lettres $A,B,C$ désignent ici des événements de lecture, non l’aire géométrique :
\begin{align}
 P_A&=\sum_{n=0}^2|j_n\rangle\langle j_n|\otimes Q_n,&
 P_B&=|0\rangle\langle0|\otimes|\chi_B\rangle\langle \chi_B|,
 \nonumber\\
 P_C&=I-P_A-P_B,& \mu(P)&=\Tr(\Sigma P).
 \label{eq:ii-tins-marks}
\end{align}
Les trois projecteurs sont orthogonaux, leur somme est l’identité et ils commutent avec
$\Sigma$. L’événement $B$ fixe une référence de rang un ; $C$ retient tout
le complément des secteurs marqués.

\Needspace{12\baselineskip}
\begin{samepage}
\begin{proposition}[Intensité relative et compression]
\label{prop:ii-tins-relative}
En écrivant $p_A=\mu(P_A)$ et $p_B=\mu(P_B)$, on obtient
\begin{equation}
 p_B=\frac{1-r}{M},\qquad
 p_A=\frac{1-r}{M}\sum_{n=0}^2d_nr^{23+3n},\qquad
 \frac{p_A}{p_B}=b_*=
 \frac{1380649}{10^{29}}=1.380649\times10^{-23}.
 \label{eq:ii-tins-relative}
\end{equation}
Si $\iota:\mathcal D\to\mathcal K$ inclut chaque secteur du porteur
dans le degré correspondant $j_n$ et dans $\operatorname{ran}Q_n$, alors
\begin{equation}
 p_B^{-1}\iota^*\Sigma\iota
 =\bigoplus_{n=0}^2r^{23+3n}I_{d_n}=\eta.
 \label{eq:ii-tins-compression}
\end{equation}
\end{proposition}
\end{samepage}
\begin{proof}
Au degré $j$, la densité agit comme $(1-r)r^j/M$ fois l’identité
de mémoire. La trace sur un sous-espace de rang $d$ multiplie ce
scalaire par $d$. L’appliquer au degré zéro et aux trois degrés marqués
donne les deux probabilités ; leur quotient annule le facteur commun.
Comme $r=10^{-1}$,
$\sum_nd_nr^{23+3n}=10^{-23}(1+380\,10^{-3}+649\,10^{-6})$,
ce qui donne la fraction indiquée. La même action scalaire, restreinte au moyen de
$\iota$, démontre \eqref{eq:ii-tins-compression}.
\end{proof}

L’intensité $b_*$ est une mesure relative, non la probabilité brute $p_A$.
Utiliser comme référence tout le vide $|0\rangle\langle0|\otimes I_M$
produirait $b_*/M$ : l’événement de référence change. Ajouter un facteur auxiliaire
indépendant et étendre les deux projecteurs par son identité conserve, en
revanche, les deux probabilités et le quotient.
% END THERMAL_R03_01

% BEGIN THERMAL_R03_02
\subsection{Appareil unitaire, récupération et domaine énergétique}
\label{subsec:ii-tins-aparato}

Dans un pointeur de dimension trois, de base $|A\rangle,|B\rangle,|C\rangle$,
définissons
\begin{equation}
 V\psi=\sum_{y=A,B,C}P_y\psi\otimes|y\rangle.
 \label{eq:ii-tins-isometry}
\end{equation}
Pour le réaliser au moyen d’un opérateur unitaire, soient $R_C=I$, $R_A$ l’échange
de $|C\rangle$ avec $|A\rangle$ et $R_B$ l’échange de $|C\rangle$
avec $|B\rangle$. Le couplage est
\begin{equation}
 \mathbb U=\sum_{y=A,B,C}P_y\otimes R_y,
 \qquad \mathbb U(\psi\otimes|C\rangle)=V\psi.
 \label{eq:ii-tins-unitary}
\end{equation}

\begin{proposition}[Enregistrement conservatif des trois alternatives]
\label{prop:ii-tins-instrument}
On a $\mathbb U^*\mathbb U=\mathbb U\mathbb U^*=I$ et $V^*V=I$.
Pour toute densité $\xi$, l’état conjoint $\widehat\xi=V\xi V^*$
permet la récupération exacte $V^*\widehat\xi V=\xi$. La lecture du
pointeur définit l’instrument
\begin{equation}
 \mathcal I_y(\xi)=P_y\xi P_y,\qquad
 p_y=\Tr(\xi P_y),\qquad
 \xi_y=\mathcal I_y(\xi)/p_y\quad(p_y>0).
 \label{eq:ii-tins-instrument}
\end{equation}
Son canal non sélectif préserve la trace et laisse $\Sigma$ invariant.
\end{proposition}
\begin{proof}
Les produits croisés des $P_y$ s’annulent. Par conséquent,
$\mathbb U^*\mathbb U=\sum_yP_y\otimes R_y^*R_y=I$, et de même
pour l’autre produit. L’orthonormalité du pointeur donne $V^*V=\sum_yP_y=I$ ;
l’appliquer aux deux membres de $\widehat\xi$ démontre la récupération.
Comprimer le pointeur sur $|y\rangle$ donne \eqref{eq:ii-tins-instrument}.
De plus, $\sum_y\Tr(P_y\xi P_y)=\Tr\xi$. Comme chaque $P_y$ commute avec
$\Sigma$, on a aussi $\sum_yP_y\Sigma P_y=\Sigma$.
\end{proof}

Pour un état présentant des cohérences entre marques, oublier le pointeur peut
les changer. La récupération précédente utilise l’état conjoint, qui les
conserve. Lire un résultat et conserver l’appareil complet sont des opérations
distinctes.

Soit $J|j\rangle=j|j\rangle$ et soit $H_{\rm sp}=f(J)\otimes I_M$, avec
$f:\mathbb N_0\to\mathbb R$. Son domaine autoadjoint est
\[
 \mathcal D(H_{\rm sp})=
 \left\{\psi:\sum_{j,a}|f(j)|^2|\psi_{j,a}|^2<\infty\right\}.
\]
Les projecteurs des marques réduisent ce domaine et commutent avec les
projections spectrales de $H_{\rm sp}$. On y vérifie
\begin{equation}
 (H_{\rm sp}\otimes I_3)V=VH_{\rm sp},\qquad
 V^*(H_{\rm sp}\otimes I_3)V=H_{\rm sp}.
 \label{eq:ii-tins-energy}
\end{equation}
En effet, appliquer $H_{\rm sp}$ à chaque terme de
\eqref{eq:ii-tins-isometry} permet de le permuter avec $P_y$ ; la norme
quadratique de la somme est la somme des normes quadratiques de ses secteurs.
La seconde égalité découle de $V^*V=I$. Les espérances
énergétiques sont ainsi conservées lorsqu’elles existent, y compris dans la réalisation
$f(j)=\epsilon_a(j+1/2)$. Un hamiltonien propre au pointeur nécessite sa
spécification : si ses trois niveaux sont dégénérés, le couplage conserve
également l’énergie totale. Cette identité n’assigne pas un coût nul à la
préparation ou à la réinitialisation physique de l’appareil.
% END THERMAL_R03_02

% BEGIN THERMAL_R03_03
\subsection{Transport tritique avec résidu et retenue}
\label{subsec:ii-tins-tritico}

La division $j=3w+b$, $0\le b<3$, définit l’opérateur unitaire
$U_3|j\rangle=|w,b\rangle$. Soit $W=U_3\otimes I_M$. Les marques transportées sont
\begin{align}
 P_A'&=\sum_{n=0}^2|7+n,2\rangle\langle7+n,2|\otimes Q_n,\nonumber\\
 P_B'&=|0,0\rangle\langle0,0|\otimes|\chi_B\rangle\langle \chi_B|,
 \qquad P_C'=I-P_A'-P_B'.
 \label{eq:ii-tins-tritic-marks}
\end{align}

\begin{proposition}[Naturalité de l’appareil]
\label{prop:ii-tins-tritic}
Si $V'$ est l’appareil des marques transportées, alors
\begin{equation}
 V'W=(W\otimes I_3)V,\qquad
 \mathcal I_y'(W\xi W^*)=W\mathcal I_y(\xi)W^*.
 \label{eq:ii-tins-square}
\end{equation}
Les probabilités et les états relatifs se transportent sans perte.
\end{proposition}
\begin{proof}
La division euclidienne donne $P_y'W=WP_y$ pour chaque $y$. Substituer ces
trois identités dans la somme définissant $V'$ prouve le premier carré ;
multiplier par les adjoints prouve le second. L’invariance de la trace
sous l’opérateur unitaire démontre l’affirmation probabiliste, et diviser par
la même probabilité positive démontre celle concernant les états relatifs.
\end{proof}

La densité transportée conserve séparément le quotient et le résidu :
\begin{equation}
 W\Sigma W^*=\rho_{r^3}\otimes\sigma_r\otimes I_M/M,
 \qquad \sigma_r=\frac{\operatorname{diag}(1,r,r^2)}{1+r+r^2}.
 \label{eq:ii-tins-tritic-state}
\end{equation}
En effet, $(1-r)r^{3w+b}=(1-r^3)(r^3)^w r^b/(1+r+r^2)$.
Pour le décalage unilatéral $S|j\rangle=|j+1\rangle$, on obtient
\begin{equation}
 U_3SU_3^*=I\otimes\bigl(|1\rangle\langle0|+|2\rangle\langle1|\bigr)
 +S_w\otimes|0\rangle\langle2|.
 \label{eq:ii-tins-carry}
\end{equation}
Les cas $b=0,1$ incrémentent le résidu ; $b=2$ le ramène à zéro et
incrémente $w$. Cette preuve conserve la retenue. Le décalage peut
changer la marque : \eqref{eq:ii-tins-square} exprime l’indépendance de l’ordre
entre observation et changement de représentation, non l’immobilité sous la dynamique.
% END THERMAL_R03_03

% BEGIN THERMAL_R03_04
\subsection{Inclusions équivalentes et extension de mémoire}
\label{subsec:ii-tins-inclusiones}

Soient $\widetilde Q_n$ d’autres projecteurs de mêmes rangs et
$\widetilde \chi_B$ un autre vecteur unitaire. À chaque degré $j_n$, il existe un opérateur unitaire
$R_{j_n}$ qui envoie $\operatorname{ran}Q_n$ sur
$\operatorname{ran}\widetilde Q_n$ : il suffit d’envoyer des bases orthonormées de ces
sous-espaces et de leurs compléments. Choisissons également $R_0\chi_B=\widetilde \chi_B$
et complétons par des opérateurs unitaires aux degrés restants. Alors
$R=\bigoplus_{j\ge0}R_j$ est unitaire, conserve $\Sigma$ et commute avec
$H_{\rm sp}$. Les identités
\begin{equation}
 \widetilde P_yR=RP_y,\qquad
 \widetilde V R=(R\otimes I_3)V
 \label{eq:ii-tins-inclusion}
\end{equation}
se vérifient bloc par bloc. Par conjugaison et cyclicité de la trace,
elles conservent les probabilités, l’énergie spectrale et $b_*$. Il n’est pas nécessaire qu’un
unique opérateur unitaire de mémoire indépendant du degré réalise simultanément
toutes les inclusions. Tout opérateur supplémentaire est transporté par le
même $R$ ; cette conjugaison n’affirme pas que sa matrice reste inchangée.

Pour une extension, soit $L:\mathbb C^M\to\mathbb C^{M'}$ une isométrie
et soit $J_L=I\otimes L$. Définissons
$P_A^+=J_LP_AJ_L^*$, $P_B^+=J_LP_BJ_L^*$ et
$P_C^+=I-P_A^+-P_B^+$. Comme $J_L^*J_L=I$, pour les trois marques,
$P_y^+J_L=J_LP_y$. Par conséquent,
\begin{equation}
 V^+J_L=(J_L\otimes I_3)V,
 \qquad \Sigma_L=J_L\Sigma J_L^*=\rho_r\otimes LL^*/M.
 \label{eq:ii-tins-extension}
\end{equation}
Le transport isométrique conserve les probabilités originales.
Si, en outre, on prépare la densité uniforme de tout l’espace ambiant étendu,
$\Sigma^+=\rho_r\otimes I_{M'}/M'$, la préparation est différente et
\begin{equation}
 p_A^+=\frac{M}{M'}p_A,\qquad
 p_B^+=\frac{M}{M'}p_B,\qquad p_A^+/p_B^+=b_*.
 \label{eq:ii-tins-uniform-extension}
\end{equation}
La preuve est de nouveau la trace du scalaire de chaque degré sur les mêmes
rangs. Dans $A$ et $B$, la normalisation annule $M/M'$ et produit les images
isométriques des états relatifs originaux ; le complément reçoit le
poids ajouté. On distingue ainsi la conservation de l’appareil, l’égalité de
l’intensité et l’égalité de l’état complet.
% END THERMAL_R03_04

% BEGIN THERMAL_R03_05
\subsection{Préparation complète, entropie et conditionnement}
\label{subsec:ii-tins-condicionamiento}

Soit $\rho$ un état d’un système d’essai fini, $H=H^*$ son hamiltonien
hérité et $\rho_0$ une préparation pure fixe. Avec $P=P_A$, définissons
\begin{equation}
 \Omega_\rho=(P\Sigma P)\otimes\rho
 +((I-P)\Sigma(I-P))\otimes\rho_0,
 \qquad \Omega_0=\Sigma\otimes\rho_0.
 \label{eq:ii-tins-preparation}
\end{equation}
Cette préparation conserve tout le complément spectral. Elle est distincte du
produit auxiliaire de \eqref{eq:ii-ter27-flag}, bien qu’elle réalise les deux mêmes
fonctionnelles sur l’essai.

\begin{theorem}[Information par référence et énergie conditionnée]
\label{thm:ii-tins-readings}
La trace de $\Omega_\rho$ vaut un, et la trace partielle sur l’espace
d’essai redonne $\Sigma$. De plus,
\begin{align}
 \mathcal S_B(\rho)&:=
 \frac{s(\Omega_\rho)-s(\Omega_0)}{p_B}=b_*s(\rho),
 \label{eq:ii-tins-information}\\
 \rho_A&:=\frac{\Tr_{\mathcal K}[(P\otimes I)\Omega_\rho(P\otimes I)]}{p_A}
 =\rho,\qquad
 \mathcal E_A(\rho)=\Tr(\rho_AH)=\Tr(\rho H).
 \label{eq:ii-tins-conditioned}
\end{align}
\end{theorem}
\begin{proof}
Les termes positifs ont pour traces $p_A$ et $1-p_A$. Comme $[P,\Sigma]=0$,
la trace partielle sur l’essai est
$P\Sigma P+(I-P)\Sigma(I-P)=\Sigma$. Soient $\lambda_i$ les valeurs propres
du bloc marqué et $p_j$ celles de $\rho$. Leurs produits remplacent
$\lambda_i$ ; l’accroissement entropique de ce bloc est
\[
 -\sum_{i,j}\lambda_ip_j\ln(\lambda_ip_j)
 +\sum_i\lambda_i\ln\lambda_i=p_As(\rho).
\]
Le complément ne change pas. Les sommes sont légitimes, car
$s(\Sigma)=-\ln(1-r)-r\ln r/(1-r)+\ln M<\infty$ et l’essai est fini.
Diviser par $p_B$ démontre \eqref{eq:ii-tins-information}. Comprimer sur
$P\otimes I$ et prendre la trace partielle donne $p_A\rho$ ; diviser par $p_A>0$
démontre \eqref{eq:ii-tins-conditioned} et son espérance énergétique.
\end{proof}

Le conditionnement conserve l’opération observationnelle du corpus. Si
$U_a$ couple l’état au registre initial $m_0$, la sortie et sa fibre sont
\[
 \operatorname{Out}_a(x)=q_a(\operatorname{pr}_M U_a(x,m_0)),\qquad
 \Omega_{a,y}=\{x:\operatorname{Out}_a(x)=y\}.
\]
Pour $\mu_a(\Omega_{a,y})>0$, on restreint la mesure à cette fibre :
\begin{equation}
 \mu_{a,y}(D)=\frac{\mu_a(D\cap\Omega_{a,y})}{\mu_a(\Omega_{a,y})}.
 \label{eq:ii-tins-fiber}
\end{equation}
Lorsque ces événements appartiennent à une même mesure spectrale projective et que
$\mu_a(D)=\Tr(\omega P_D)$, on a
$P_DP_y=P_{D\cap\Omega_{a,y}}$. La cyclicité de la trace prouve
\[
 \Tr\!\left(\frac{P_y\omega P_y}{\Tr(\omega P_y)}P_D\right)
 =\mu_{a,y}(D).
\]
L’égalité s’étend aux observables bornées de cette algèbre par sommes
spectrales et approximation uniforme. Appliquée à l’événement $A$ de
\eqref{eq:ii-tins-preparation}, elle donne exactement l’énergie moyenne
conditionnée précédente. Elle ne nécessite pas d’hypothèse dynamique d’ergodicité.
Une purification supplémentaire, lorsqu’elle existe, est conservée dans son propre facteur ;
l’entropie calculée ici appartient à l’algèbre observable spécifiée.
% END THERMAL_R03_05

% BEGIN THERMAL_R03_06
\subsection{Caractérisation par additivité et raffinement}
\label{subsec:ii-tins-aditividad}

L’évaluation précédente admet une caractérisation qui fixe son poids sans
prendre $b_*$ comme coefficient postulé. L’état d’essai $\rho$ étant fixé,
considérons une contribution $F_\rho(x)$ pour une alternative de poids
reçu $x\in[0,1]$. Dans la lecture relative, ce poids se mesure par rapport à
la référence $B$ ; seules les alternatives du domaine indiqué sont considérées.
Nous exigeons : la dépendance à l’égard du poids et non du nom de l’alternative ; l’additivité
$F_\rho(x+y)=F_\rho(x)+F_\rho(y)$ pour $x+y\le1$ ; la positivité ; et la
normalisation $F_\rho(1)=s(\rho)$. Le domaine doit réaliser ces répartitions,
ou les raffinements ordonnés décrits ci-après.

\begin{theorem}[Unicité du poids informationnel additif]
\label{thm:ii-tins-additive}
Sur le domaine scalaire complet, les quatre conditions précédentes impliquent
$F_\rho(x)=xs(\rho)$. La même conclusion vaut pour une réalisation
d’alternatives avec des raffinements équipondérés de poids $3^{-n}$ et, pour
chaque marque $P$, des approximations $P_n^-\le P\le P_n^+$ dont les poids convergent
vers celui de $P$ et dont les contributions appartiennent à ces raffinements.
\end{theorem}
\begin{proof}
L’additivité donne $F_\rho(0)=0$. Si $x\le y$, alors
$F_\rho(y)-F_\rho(x)=F_\rho(y-x)\ge0$, de sorte que le lecteur est monotone.
La partition de l’unité en $3^n$ alternatives égales donne
$F_\rho(3^{-n})=3^{-n}s(\rho)$, et en réunir $k$ donne
$F_\rho(k3^{-n})=k3^{-n}s(\rho)$. Pour $0\le x<1$, nous prenons
\[
 \ell_n=3^{-n}\lfloor3^nx\rfloor,\qquad u_n=\ell_n+3^{-n}.
\]
Ce sont des bornes de $x$ dans $[0,1]$. La monotonie implique
$\ell_ns(\rho)\le F_\rho(x)\le u_ns(\rho)$ ; leur différence
$3^{-n}s(\rho)$ tend vers zéro. Cela prouve la formule, y compris
$s(\rho)=0$ ; $x=1$ est la normalisation. Pour des marques ordonnées, on répète
l’argument avec $P_n^-$ et $P_n^+$, car l’additivité positive implique
l’isotonie. On ne présuppose pas la continuité différentiable du lecteur.
\end{proof}

En particulier, la contribution de $A$ est $b_*s(\rho)$ lorsque sa
mesure relative est conservée. La caractérisation s’applique également à la marque
normalisée de \eqref{eq:ii-ter27-flag}. La fonctionnelle
$-\Tr[(P\otimes I)(\widehat\eta\otimes\rho)(I\otimes\ln\rho)]$
y satisfait les conditions par linéarité de la trace et assigne $s(\rho)$
à la marque certaine.

Le raffinement nécessite une réalisation effective : le projecteur $P_B$ de
rang un ne contient pas trois sous-projecteurs non nuls équipondérés. Une
extension auxiliaire peut réaliser des raffinements, mais doit transporter
la mesure et la lecture. De même, trois descendants d’un cylindre
n’impliquent pas, par leur cardinalité, des poids égaux. Pour un poids rationnel, il suffit
d’une partition équipondérée du dénominateur correspondant, si elle existe dans
la réalisation ; la preuve ternaire ne présuppose pas ce dénominateur particulier.

Les hypothèses séparent exactement les libertés écartées. Le lecteur
$f(x)=x^2$ respecte les produits mais donne
$f(x+y)-f(x)-f(y)=2xy$ ; avec $x=1/3$, $y=1/5$, le défaut est $2/15$.
Sans normalisation, $f(x)=cx$ conserve la positivité et l’additivité. Sans
identifier la mesure reçue, $(1/3,2/3)$ et $(1/2,1/2)$ sont deux mesures
positives, additives et normalisées différentes sur deux alternatives.
L’entropie fixe du registre des étiquettes reste à part : le théorème
n’élimine pas le terme de mélange d’une entropie conjointe.
% END THERMAL_R03_06

% BEGIN THERMAL_R03_07
\subsection{Réalisation par visites entre retours}
\label{subsec:ii-tins-retornos}

L’intensité relative possède également une interprétation dynamique précise.
Soit $(X,\mu,T)$ un système probabiliste inversible, ergodique et préservant
la mesure, avec des événements $A,B$ de mesures $p_A,p_B>0$. Pour $x\in B$, soient
\[
 R_B(x)=\min\{k\ge1:T^kx\in B\},\qquad
 C_A(x)=\sum_{k=0}^{R_B(x)-1}\mathbf1_A(T^kx).
\]
L’excursion conserve sa séquence complète ; $C_A$ n’est qu’une de ses observables.

\begin{theorem}[Récompense par retour à une référence]
\label{thm:ii-tins-return}
Avec $\mu_B=\mu|_B/p_B$, pour tout $f\in L^1(\mu)$,
\begin{equation}
 \mathbb E_{\mu_B}\!\left[\sum_{k=0}^{R_B-1}f\circ T^k\right]
 =\frac1{p_B}\int_Xf\,d\mu.
 \label{eq:ii-tins-return}
\end{equation}
En particulier, $\mathbb E_{\mu_B}C_A=p_A/p_B=b_*$.
\end{theorem}
\begin{proof}
Soient $B_k=\{x\in B:R_B(x)>k\}$. La récurrence, l’inversibilité et
l’ergodicité assurent que les niveaux $T^kB_k$, $k\ge0$, partitionnent $X$
à un ensemble négligeable près : chaque point a une dernière visite à $B$ dans son
passé. Pour $f\ge0$, la préservation de la mesure et la convergence monotone donnent
\[
 \int_B\sum_{k=0}^{R_B-1}f(T^kx)\,d\mu
 =\sum_{k\ge0}\int_{T^kB_k}f\,d\mu=\int_Xf\,d\mu.
\]
Pour $f\in L^1$, appliquer d’abord l’identité à $|f|$ justifie
l’intégrabilité, puis on soustrait les parties positive et négative.
Diviser par $p_B$ prouve la formule ; $f=\mathbf1_A$ donne la conséquence.
\end{proof}

Dans une chaîne finie irréductible de matrice $P$ et de distribution stationnaire
$\pi>0$, la preuve admet une forme algébrique. Partant de l’état de
référence $B$, soit $v_i$ le nombre espéré de visites à $i$ avant le
premier retour positif. Le retour a une espérance finie ; sommer les
transitions jusqu’à celui-ci montre que les arrivées et les départs ont les mêmes
espérances, donc $vP=v$. Comme $v_B=1$, l’unicité de la distribution
stationnaire donne $v_i=\pi_i/\pi_B$. Sommer sur $A$ reproduit
\eqref{eq:ii-tins-return}. Regrouper la distribution spectrale en $B$, les
trois secteurs de $A$ et le complément fournit une réalisation finie
avec le quotient exact \eqref{eq:ii-tins-relative}.

Si chaque visite à $A$ porte un essai $\rho$, indépendant des autres
conditionnellement au registre de l’excursion, les essais forment
$\rho^{\otimes C_A}$. L’additivité tensorielle donne
\begin{align}
 \mathbb E_{\mu_B}s(\rho^{\otimes C_A})&=b_*s(\rho),\nonumber\\
 \frac{\mathbb E_{\mu_B}[C_A\Tr(\rho H)]}
      {\mathbb E_{\mu_B}C_A}&=\Tr(\rho H).
 \label{eq:ii-tins-visit-readings}
\end{align}
La seconde expression est le quotient de l’énergie totale par le nombre total de visites,
non une moyenne de quotients par excursion, qui seraient indéfinis lorsque
$C_A=0$. L’énergie accumulée par retour vaut $b_*\Tr(\rho H)$.

L’information par retour et l’énergie par visite réalisent donc
les deux dénominateurs de \eqref{eq:ii-tins-information}--
\eqref{eq:ii-tins-conditioned}. Si les essais sont corrélés,
on utilise leur entropie conditionnelle conjointe. Un processus non inversible admet
la formulation dans son extension bilatérale stationnaire ; sans ergodicité,
le domaine doit se restreindre au support balayé par la section de retour.
Ces hypothèses dynamiques ne sont pas nécessaires au théorème opératoriel
\ref{thm:ii-tins-readings}. Ni une chaîne stationnaire d’essai ne s’identifie
à elle seule à la mise à jour TPK, ni $R_B$ à la période matérielle
$108t_0$.
% END THERMAL_R03_07

% BEGIN THERMAL_R03_08
\subsection{Référence commune et corrélations entre systèmes}
\label{subsec:ii-tins-comun}

Une seule référence et $m$ systèmes indépendants produisent le registre
$\widehat\eta\otimes\rho_1\otimes\cdots\otimes\rho_m$, avec une
unique marque $P\otimes I$. Son poids reste $b_*$. Pour
$H_{\rm tot}=\sum_jH_j$, avec les extensions tensorielles correspondantes,
\begin{equation}
 \mathcal S_{P,\rm tot}=b_*\sum_js(\rho_j),\qquad
 \mathcal E_{P,\rm tot}=\sum_j\Tr(\rho_jH_j).
 \label{eq:ii-tins-common}
\end{equation}
La preuve est la factorisation de la trace et
$\ln(\rho_1\otimes\rho_2)=\ln\rho_1\otimes I+I\otimes\ln\rho_2$
sur le support, itérée sur les facteurs. Pour un état corrélé
$\rho_{1\ldots m}$, la même marque donne $b_*s(\rho_{1\ldots m})$. Dans deux
systèmes,
\begin{equation}
 \mathcal S_{P,1}+\mathcal S_{P,2}-\mathcal S_{P,12}
 =b_*I(1:2)_\rho,
 \label{eq:ii-tins-mutual}
\end{equation}
par la définition de l’information mutuelle. Un hamiltonien avec interaction
conserve également une seule référence, bien que son état de Gibbs ne se factorise pas.

Deux références indépendantes de marque conjointe $P_1\otimes P_2$
ont pour poids $b_1b_2$ ; pour des références corrélées, le poids est
$\Tr[\omega_{R_1R_2}(P_1\otimes P_2)]$ et ne se factorise pas nécessairement.
Copier une étiquette ne crée pas d’indépendance. En effet, l’isométrie
\[
 V_Pv=Pv\otimes|1\rangle+(I-P)v\otimes|0\rangle
\]
satisfait $V_P^*(P\otimes|1\rangle\langle1|)V_P=P$ ; la trace sur
l’état transporté conserve $b_*$, non $b_*^2$. Plus généralement,
pour une isométrie $L$,
$\Tr[(L\omega L^*)(LPL^*)]=\Tr(\omega P)$, par $L^*L=I$.
Revenir en conservant le même registre n’élève pas automatiquement son poids
à une puissance. Une nouvelle sélection de survie utilise sa propre marque.
% END THERMAL_R03_08

% BEGIN THERMAL_R03_09
\subsection{Sélection thermométrique et composition avec l’action et l’horizon}
\label{subsec:ii-tins-seleccion}

Fixons la référence et un hamiltonien fini non scalaire $H$. Pour
$\rho_\beta=e^{-\beta H}/Z(\beta)$, $\beta>0$, le théorème additif et
l’énergie conditionnée produisent la paire
$\mathcal S_B=b_*s(\rho_\beta)$, $\mathcal E_A=\Tr(\rho_\beta H)$.
Si $E=\Tr(\rho_\beta H)$, la somme spectrale finie donne
$E'=-\operatorname{Var}_{\rho_\beta}(H)<0$ et
$(\ln Z)'=-E$. Dériver $s=\beta E+\ln Z$ donne $s'=\beta E'$ et, par conséquent,
\begin{equation}
 \frac{d\mathcal S_B}{d\mathcal E_A}=b_*\beta,\qquad
 \Theta_{B,A}=\frac1{b_*\beta},\qquad B_*(u_\Theta)=b_*u_E.
 \label{eq:ii-tins-temperature}
\end{equation}
Les bases positives satisfont $u_S=u_E/u_\Theta$. La dernière équation
détermine l’unique application linéaire positive qui envoie
$\Theta_{B,A}u_\Theta$ sur $\beta^{-1}u_E$. Le bilan
\[
 \Tr(\rho H)-b_*\Theta s(\rho)-
 \bigl[\Tr(\rho_\Theta^*H)-b_*\Theta s(\rho_\Theta^*)\bigr]
 =b_*\Theta D(\rho\Vert\rho_\Theta^*)
\]
avec $\rho_\Theta^*\propto e^{-H/(b_*\Theta)}$ conserve le minimum
unique démontré dans \ref{cor:ii-ter27-variational}. La sélection additive
caractérise maintenant le poids qui entre dans ce bilan.

Avec $\epsilon_a=h_a/(108t_0)=\hbar_ac/L_\alpha$,
$q_A=\gamma a_0L_\alpha^2$ et
$\mathcal T_{I/A}=\epsilon_a/q_A$, les sections thermiques sont
\begin{align}
 b_*\Theta_{{\rm clk},P,a}
 &=\frac{h_a}{108t_0\ln3}
 =\frac{\mathcal T_{I/A}q_A}{\ln3},\nonumber\\
 b_*\Theta_{H,P}(R)&=\frac{\hbar_ac}{2\pi R},&
 \frac{\Theta_{H,P}(R)}{\Theta_{{\rm clk},P,a}}
 &=\frac{L_\alpha}{R}\frac{\ln3}{2\pi}.
 \label{eq:ii-tins-action-area}
\end{align}
Ce sont des équations de coordonnées de $B_*$ dans les bases transportées.
La première utilise la section $\beta_{{\rm clk},a}=\ln3/\epsilon_a$ ;
la seconde, l’échelle $\tau_R=\hbar_ac/(2\pi R)$. Les diviser prouve
le quotient. La relation radiale $G_a\hbar_a=c^3L_\alpha^2$ implique
\begin{equation}
 G_ab_*\Theta_{{\rm clk},P,a}\ln3=c^4L_\alpha.
 \label{eq:ii-tins-gravity}
\end{equation}
À l’horizon, $\mathcal S_B=b_*s$ et
$\Theta_{H,P}=\tau_R/b_*$ convertissent le bilan précédent en
\[
 \delta E_R-\Theta_{H,P}\delta\mathcal S_B
 =\tau_RD(\sigma\Vert\rho_A).
\]
La composition maintient la résolution de Barbero, l’incidence sectorielle,
l’action et l’origine énergétique : l’annulation des échelles conserve la
même énergie, elle n’ajoute pas de paramètre métrologique au constructeur.

La sélection a un domaine défini. Si l’on conditionne les deux canaux,
on obtient $(E,s)$ ; si l’on compte les deux par référence, $(b_*E,b_*s)$.
Dans les deux cas, la dérivée est $\beta$. La paire caractérisée ici est
$(E,b_*s)$ et utilise l’énergie par visite et l’information par référence.
La conservation de l’archive permet les trois observations, mais ne les
identifie pas. Les produits énergétiques ne sélectionnent pas non plus à eux seuls le poids :
pour $p>0$, remplacer $b_*$ par $b_*^p$ et diviser les températures par
le même facteur relatif conserve ces produits et la composition en série ;
le lecteur quadratique ne satisfait cependant pas l’additivité démontrée ci-dessus.

L’instrument des marques, sa récupération et son
transport, la contribution informationnelle additive, l’énergie conditionnée
et le transducteur de cette réalisation sont déterminés. L’identifier à un instrument
thermométrique du constructeur complet exige de transporter vers $A$ et $B$
ses événements, sa mesure et sa paire de normalisations, en conservant le
raffinement et la mémoire. L’existence de l’appareil ne remplace pas cette
identification particulière. Si la référence change avec $\beta$, ses dérivées
apparaissent en outre ; si la projection ne commute pas avec $\Sigma$, la préparation
\eqref{eq:ii-tins-preparation} conserve sa diagonale par blocs et pas
nécessairement la densité originale. Ces conditions précisent la portée
des égalités, sans altérer les résultats antérieurs d’action, d’aire,
de capacité et de transport.
% END THERMAL_R03_09
